\begin{lemma}[The family of jets is compact]
Let $1<p<2$ and let $(\kappa,C_0)$ be a first-order estimate datum for $p$.  Then the family
$\mathcal{K}_p$ is a compact subset of $X_{1/2}\to^{b}(\mathbb{R}\times\mathrm{Plane})$.
\end{lemma}
\begin{proof}
Write $\mathcal{G}_p$ for the generating set, so that $\mathcal{K}_p=\overline{\mathcal{G}_p}$ by
\noderef{CompactJetFamily}.  Throughout, $\kappa$ and $C_0$ are the data of the first-order
estimate datum, so $0<\kappa<1$ and $C_0>0$.

\medskip\noindent\emph{Step 0: two inclusions of balls.}
Let $x\in X_{1/2}$, so $\lvert x\rvert\le\tfrac12$ as a point of the plane.  Then
$\lvert x\rvert\le\tfrac12<\tfrac34$ and $\lvert x\rvert<1$, so
\begin{equation}\label{eq:cjfc-incl}
  X_{1/2}\subseteq B_{3/4}\subseteq B_1 ,
\end{equation}
the second inclusion because $\tfrac34<1$.

\medskip\noindent\emph{Step 1: the estimates available on a generator.}
Let $z\in\mathcal{G}_p$.  By \noderef{CompactJetFamily} there are
$v\colon\mathrm{Plane}\to\mathbb{R}$ and $g\colon\mathrm{Plane}\to\mathrm{Plane}$ with $v$
$p$-harmonic in $B_1$ with weak gradient $g$, with $\lvert v(y)\rvert\le1$ for $y\in B_1$, with
$v(0)=0$, and with $z(x)=\bigl(v(x),Dv(x)\bigr)$ for every $x\in X_{1/2}$.

Apply \noderef{PHarmJetEquicontinuity} with $x_0:=0$, $R:=\tfrac12$ and $M:=1$.  Its hypotheses
hold: $1<p<2$; $R=\tfrac12>0$ and $M=1>0$; $v$ is $p$-harmonic in $B_{2R}(0)=B_1$ with weak
gradient $g$, since $2\cdot\tfrac12=1$; and $\lvert v(y)\rvert\le1=M$ for every
$y\in B_{2R}(0)=B_1$.  Its conclusion, evaluated at these parameters --- where
$3R/2=\tfrac34$, $C_0M/(2R)=C_0\cdot1/1=C_0$ and $2R=1$ --- is that $v$ is of class $C^1$ on
$B_{3/4}$ and
\begin{equation}\label{eq:cjfc-grad}
  \lvert Dv(y)\rvert\le C_0\qquad(y\in B_{3/4}),
\end{equation}
\begin{equation}\label{eq:cjfc-hol}
  \lvert Dv(y)-Dv(y')\rvert\le C_0\,\lvert y-y'\rvert^{\kappa}\qquad(y,y'\in B_{3/4}).
\end{equation}

\medskip\noindent\emph{Step 2: a compact target set containing every jet value.}
Put
\[
  s:=[-1,1]\times\overline{B_{C_0}(0)}\subseteq\mathbb{R}\times\mathrm{Plane},
\]
which is compact by clause (2) of the background proposition on bounded continuous functions,
applied with $A:=1$ and $B:=C_0$.

We claim $z(x)\in s$ for every $z\in\mathcal{G}_p$ and every $x\in X_{1/2}$.  Take $v,g$ as in
Step 1.  The first component of $z(x)$ is $v(x)$, and $x\in B_1$ by \eqref{eq:cjfc-incl}, so
$\lvert v(x)\rvert\le1$, that is $v(x)\in[-1,1]$.  The second component is $Dv(x)$, and
$x\in B_{3/4}$ by \eqref{eq:cjfc-incl}, so $\lvert Dv(x)\rvert\le C_0$ by \eqref{eq:cjfc-grad},
that is $Dv(x)\in\overline{B_{C_0}(0)}$.  Hence $z(x)\in s$.

\medskip\noindent\emph{Step 3: equicontinuity, with a modulus uniform over the family.}
Let $\varepsilon>0$.  Since $\kappa>0$, clause (7) of the background proposition on elementary
properties of real powers gives that $t\mapsto t^{\kappa}$ is continuous on $[0,\infty)$ with
value $0$ at $t=0$, so there is $s_0>0$ with $C_0\,s_0^{\kappa}\le\varepsilon$.  Put
\[
  \delta:=\min\Bigl\{s_0,\ \frac{\varepsilon}{C_0}\Bigr\}>0 ,
\]
positive because $s_0>0$, $\varepsilon>0$ and $C_0>0$.  Let $z\in\mathcal{G}_p$ with witnesses
$v,g$ as in Step 1, and let $x,x'\in X_{1/2}$ with $\lvert x-x'\rvert<\delta$.

\emph{First components.}  The ball $B_{3/4}$ is convex; $v$ is of class $C^1$ on $B_{3/4}$ by
Step 1, hence differentiable at every point of $B_{3/4}$ by clause (2) of the background
proposition on locality of $C^1$ regularity; and $\lvert Dv\rvert\le C_0$ on $B_{3/4}$ by
\eqref{eq:cjfc-grad}.  Both $x$ and $x'$ lie in $B_{3/4}$ by \eqref{eq:cjfc-incl}.  The
background mean value inequality on a convex set therefore gives
\[
  \lvert v(x)-v(x')\rvert\le C_0\,\lvert x-x'\rvert<C_0\,\delta
  \le C_0\cdot\frac{\varepsilon}{C_0}=\varepsilon .
\]

\emph{Second components.}  By \eqref{eq:cjfc-hol} with $y:=x$, $y':=x'$, both in $B_{3/4}$, and
since $\lvert x-x'\rvert<\delta\le s_0$ with $t\mapsto t^{\kappa}$ nondecreasing on $[0,\infty)$
(clause (3) of the elementary-powers proposition),
\[
  \lvert Dv(x)-Dv(x')\rvert\le C_0\,\lvert x-x'\rvert^{\kappa}\le C_0\,s_0^{\kappa}\le\varepsilon .
\]

By clause (2) of the bounded-continuous-functions proposition the distance in
$\mathbb{R}\times\mathrm{Plane}$ is the maximum of the two component distances, so the two
displays give
\[
  d\bigl(z(x),z(x')\bigr)
  =\max\bigl(\lvert v(x)-v(x')\rvert,\ \lvert Dv(x)-Dv(x')\rvert\bigr)\le\varepsilon .
\]
The number $\delta$ was chosen from $\varepsilon$, $\kappa$ and $C_0$ alone, and $\kappa$ and
$C_0$ are fixed with the datum before any $z$ is chosen; so this is a modulus uniform over
$\mathcal{G}_p$.  By clause (4) of the bounded-continuous-functions proposition --- a uniform
modulus suffices for equicontinuity --- the family of underlying maps of the elements of
$\mathcal{G}_p$ is equicontinuous.

\medskip\noindent\emph{Step 4: Arzel\`a--Ascoli.}
The base space $X_{1/2}$ is compact, as recorded in \noderef{XHalf}.  Clause (5) of the
bounded-continuous-functions background proposition, applied with $K:=X_{1/2}$,
$Y:=\mathbb{R}\times\mathrm{Plane}$, the compact set $s$ of Step 2 and the set
$A:=\mathcal{G}_p$, has exactly two further hypotheses: that $f(x)\in s$ for every
$f\in\mathcal{G}_p$ and every $x\in X_{1/2}$, which is Step 2, and that the family of underlying
maps of the elements of $\mathcal{G}_p$ be equicontinuous, which is Step 3.  Its conclusion is
that $\overline{\mathcal{G}_p}$ is compact.  By \noderef{CompactJetFamily},
$\overline{\mathcal{G}_p}=\mathcal{K}_p$, so $\mathcal{K}_p$ is compact.  Note that no
closedness of $\mathcal{G}_p$ was needed: the compactness assertion is about the closure, which
is what the paper's family is.
\end{proof}
